\documentclass[10pt]{article}
\usepackage[a4paper,margin=0.86in]{geometry}
\renewcommand{\rmdefault}{cmr}
\usepackage{amsmath,amssymb,amsthm,mathtools,amscd,enumitem,xcolor,hyperref}
\hypersetup{colorlinks=true,linkcolor=blue!55!black,citecolor=blue!55!black,urlcolor=blue!55!black}
\setlength{\emergencystretch}{4em}
\setlength{\parindent}{0pt}
\setlength{\parskip}{0.52em}
\setlist{leftmargin=1.55em,itemsep=0.2em,topsep=0.25em}
\sloppy

\newtheorem{theorem}{Theorem}[section]
\newtheorem{proposition}[theorem]{Proposition}
\newtheorem{lemma}[theorem]{Lemma}
\newtheorem{corollary}[theorem]{Corollary}
\theoremstyle{definition}
\newtheorem{definition}[theorem]{Definition}
\newtheorem{convention}[theorem]{Convention}
\newtheorem{remark}[theorem]{Remark}

\newcommand{\Zl}{\mathbf Z_\ell}
\newcommand{\Ql}{\mathbf Q_\ell}
\newcommand{\Qlb}{\overline{\mathbf Q}_\ell}
\newcommand{\OO}{\mathcal O}
\newcommand{\Dbc}{D^b_c}
\newcommand{\Dctf}{D^b_{\mathrm{ctf}}}
\newcommand{\red}{\operatorname{red}}
\newcommand{\Tr}{\operatorname{Tr}}

\begin{document}
\begin{center}
{\large\bfseries Smooth Purity and Its Trace Normalization\par}
\medskip
{\normalsize Orientation classes, purity, and exceptional pullback\par}
\end{center}

\begin{abstract}
For a smooth separated morphism \(f:X\to Y\) of pure relative dimension
\(d\), smooth purity identifies \(Rf^!\) with the shifted and twisted inverse
image \(Lf^*(d)[2d]\).  The statement is decomposed here into the Tate-line
construction, the normalized affine-line class, the affine-space product,
etale localization, gluing, coherence, and coefficient transport.  The
resulting trace is fixed by a geometric-fiber normalization.
\end{abstract}

\tableofcontents

\section{Coefficient and categorical conventions}

Fix a prime \(\ell\).  Schemes are noetherian of finite Krull dimension and
\(\ell\) is invertible on them.  The topology is the small pro-etale
topology, with the finite-coefficient comparison to the small etale topology.

\begin{definition}[Finite coefficient categories]
Let \(\Lambda\) be a finite local \(\mathbf Z/\ell^n\mathbf Z\)-algebra.  Put

\[
 D(X,\Lambda)=D\bigl(\mathsf{Mod}(X_{\mathrm{pro\acute et}},\Lambda_X)\bigr).
\]

The full subcategory \(\Dbc(X,\Lambda)\) consists of complexes with bounded
constructible cohomology.  The subcategory \(\Dctf(X,\Lambda)\) consists of
objects of \(\Dbc\) having finite Tor-amplitude over \(\Lambda_X\).  All are
ordinary Verdier-localized triangulated categories.
\end{definition}

\begin{definition}[Adic, rational, and algebraically closed stages]
Let \(E/\Ql\) be finite, let \(\OO_E\) be its ring of integers, choose a
uniformizer \(\varpi\), and put \(\Lambda_n=\OO_E/\varpi^n\).  The category
\(\Dbc(X,\OO_E)\) is the triangulated category of normalized systems
\(K_\bullet=(K_n)\), with \(K_n\in\Dbc(X,\Lambda_n)\), coherent isomorphisms

\[
 K_{n+1}\otimes^L_{\Lambda_{n+1,X}}\Lambda_{n,X}\xrightarrow{\sim}K_n,
\]

and common bounds and stratifications.  Define \(\Dctf(X,\OO_E)\) by a
common finite Tor-amplitude bound.  The \(E\)-categories are the
idempotent-complete essential images after inverting \(\varpi\).  The
\(\Qlb\)-categories are the Karoubi completions of the filtered \(2\)-colimits
over finite extensions \(E/\Ql\).
\end{definition}

\begin{definition}[Finite reduction]
For \(m\leq n\), derived extension of scalars is the exact functor

\[
 \red_{n,m,X}:D(X,\Lambda_n)\longrightarrow D(X,\Lambda_m),
 \qquad
 K\longmapsto K\otimes^L_{\Lambda_{n,X}}\Lambda_{m,X}.
\]

It preserves \(\Dctf\), and the canonical maps
\(\red_{m,r}\red_{n,m}\cong\red_{n,r}\) satisfy the cocycle equation.
\end{definition}

\begin{convention}[Finite-stage scope]
From the affine-line orientation through the finite-stage smooth dualizing
complex, the coefficient ring is a fixed finite local
\(\mathbf Z/\ell^n\mathbf Z\)-algebra \(\Lambda\).  Statements for
\(\OO_E\), \(E\), and \(\Qlb\) occur only after the finite-reduction theorem
and the normalized-system reconstruction theorem in the coefficient-transport
section.
\end{convention}

\begin{definition}[Ambient compact-support functor]
At a finite coefficient stage, choose a compactification
\(X\xrightarrow j\overline X\xrightarrow pY\) of a separated finite-type
morphism \(f:X\to Y\).  Define

\[
 Rf_!=Rp_*j_!:D(X,\Lambda)\longrightarrow D(Y,\Lambda),
\]

where \(j_!\) is exact extension by zero and \(Rp_*\) is computed on a
K-injective replacement.  Proper comparison along common refinements makes
this independent of the compactification.
\end{definition}

\begin{theorem}[Ambient exceptional adjoint]
Assume the small pro-etale site is compactly generated by the shifts of
affine quasi-compact weakly contractible objects, and assume that proper
direct image for the chosen compactification has a uniform cohomological
dimension on arbitrary pro-etale \(\Lambda\)-module sheaves.  Then \(Rf_!\)
preserves arbitrary coproducts.  Brown representability in the ordinary
triangulated category \(D(X,\Lambda)\) gives an exact right adjoint

\[
 Rf^!:D(Y,\Lambda)\longrightarrow D(X,\Lambda)
\]

and natural bijections

\[
 \operatorname{Hom}_Y(Rf_!K,L)
 \cong\operatorname{Hom}_X(K,Rf^!L).
\]

Its counit is the exceptional trace

\[
 \Tr_f:Rf_!Rf^!\longrightarrow\operatorname{id}_{D(Y,\Lambda)}.
\]
\end{theorem}

\begin{theorem}[Bounded constructibility and finite-Tor preservation]
Let \(f\) satisfy the ambient hypotheses above.  The finite compact-support
amplitude theorem and constructible devissage imply

\[
 Rf^!: \Dbc(Y,\Lambda)\longrightarrow\Dbc(X,\Lambda).
\]

The stratified finite-Tor theorem further implies

\[
 Rf^!: \Dctf(Y,\Lambda)\longrightarrow\Dctf(X,\Lambda).
\]

Both are restrictions of the ambient right adjoint, and \(Rf_!\dashv Rf^!\)
on each displayed pair of subcategories.  Finite reduction, normalized adic
reconstruction, scalar extension, and idempotent completion give these
restricted adjunctions for \(A=\OO_E,E,\Qlb\).
\end{theorem}

\begin{definition}[Right-adjoint base-change mate]
For a Cartesian square

\[
\begin{CD}
X' @>{u}>> X\\
@V{f'}VV @VV{f}V\\
Y' @>{v}>> Y,
\end{CD}
\]

compact-support base change is

\[
 \operatorname{BC}_!:Lv^*Rf_!\xrightarrow{\sim}Rf'_!Lu^*.
\]

Its right-adjoint mate is

\[
 Ru_*Rf'^!\longrightarrow Rf^!Rv_*.
\]

The inverse mate is denoted

\[
 \operatorname{BC}^!_*:Rf^!Rv_*\xrightarrow{\sim}Ru_*Rf'^!.
\]
\end{definition}

\begin{theorem}[Composition and etale normalization of \(Rf^!\)]
For separated finite-type maps \(X\xrightarrow fY\xrightarrow gZ\), there is
a compact-support composition isomorphism

\[
 \alpha_{g,f}:Rg_!Rf_!\xrightarrow{\sim}R(gf)_!,
\]

and its right-adjoint mate is the coherent isomorphism

\[
 c_{g,f}:R(gf)^!\xrightarrow{\sim}Rf^!Rg^!.
\]

If \(u:U\to X\) is etale, then

\[
 Ru^!\xrightarrow{\sim}Lu^*:D(X,A)\longrightarrow D(U,A).
\]

For three composable maps, the two composites of the \(\alpha\)'s agree, and
the two composites of the \(c\)'s agree.  These identifications satisfy the
identity laws and etale base-change coherence and are compatible with the
counits.
\end{theorem}

\section{Tate lines and twists}

\begin{definition}[Finite Tate line]
For \(\Lambda=\mathbf Z/\ell^n\mathbf Z\), define the finite Tate line by

\[
 \Lambda_X(1)=\mu_{\ell^n,X}
 =\ker\bigl([\ell^n]:\mathbf G_{m,X}\to\mathbf G_{m,X}\bigr).
\]

It is a locally free rank-one \(\Lambda_X\)-module sheaf.  For a finite
\(\mathbf Z/\ell^n\mathbf Z\)-algebra \(\Lambda'\), put

\[
 \Lambda'_X(1)=\Lambda_X(1)\otimes_{\mathbf Z/\ell^n\mathbf Z}\Lambda'.
\]
\end{definition}

\begin{definition}[Adic and rational Tate lines]
The transition map
\(\mu_{\ell^{n+1},X}\to\mu_{\ell^n,X}\) is
\(\zeta\mapsto\zeta^\ell\).  The normalized system of these sheaves defines

\[
 (\Zl)_X(1)=(\mu_{\ell^n,X})_{n\geq1}.
\]

The \((\OO_E)_X(1)\), \(E_X(1)\), and \((\Qlb)_X(1)\) Tate lines are obtained
by scalar extension, filtered \(2\)-colimit, and Karoubi completion.
\end{definition}

\begin{definition}[Derived internal Hom]
At a finite coefficient stage \(A=\Lambda\), for cochain complexes \(P,I\)
of \(A_X\)-module sheaves, put

\[
 \mathcal Hom_X^n(P,I)=
 \prod_{a\in\mathbf Z}\mathcal Hom_{A_X}(P^a,I^{a+n}),
 \qquad
 d(\psi)=d_I\psi-(-1)^n\psi d_P.
\]

Choose a functorial K-flat replacement \(P(M)\to M\) and a functorial
K-injective replacement \(N\to I(N)\).  Define

\[
 R\mathcal Hom_X(M,N)=
 \mathcal Hom_X^\bullet(P(M),I(N)).
\]

Levelwise reduction, normalized-system reconstruction, scalar extension,
and idempotent completion define the same bifunctor at the
\(\OO_E\), \(E\), and \(\Qlb\) stages.
\end{definition}

\begin{definition}[Twist]
For any coefficient stage \(A\), the Tate line \(A_X(1)\) is invertible for
the derived tensor product.  Define

\[
 A_X(r)=
 \begin{cases}
 A_X(1)^{\otimes r},&r\geq0,\\
 R\mathcal Hom_X(A_X(-r),A_X),&r<0,
 \end{cases}
 \qquad
 K(r)=K\otimes^L_{A_X}A_X(r).
\]

Evaluation and coevaluation make \(A_X(-1)\) the tensor inverse of
\(A_X(1)\).
\end{definition}

\begin{proposition}[Additivity and pullback of twists]
For all \(r,s\in\mathbf Z\), there are coherent isomorphisms

\[
 (K(r))(s)\xrightarrow{\sim}K(r+s),
 \qquad
 Lf^*(K(r))\xrightarrow{\sim}(Lf^*K)(r).
\]

They satisfy the associativity pentagon, the zero-twist unit laws, symmetry,
and compatibility with reduction and scalar extension.
\end{proposition}

\begin{proposition}[Frobenius normalization]
If \(X\) descends to \(\mathbf F_q\) and geometric Frobenius is used, its
action on \(A_X(1)\) is multiplication by \(q^{-1}\), and its action on
\(A_X(r)\) is \(q^{-r}\).
\end{proposition}

\section{The normalized affine-line class}

Let \(p:\mathbf A^1_Y\to Y\) be the projection, let
\(s:Y\hookrightarrow\mathbf A^1_Y\) be the zero section, and let
\(j:\mathbf G_{m,Y}\hookrightarrow\mathbf A^1_Y\) be its open complement.

\begin{definition}[Cohomology with support]
For a closed immersion \(i:Z\hookrightarrow X\) with open complement
\(j:U\hookrightarrow X\), define

\[
 R\Gamma_Z(K)=i_*Ri^!K.
\]

It is characterized by the localization distinguished triangle

\[
 i_*Ri^!K\longrightarrow K\longrightarrow Rj_*j^*K
 \longrightarrow i_*Ri^!K[1].
\]
\end{definition}

\begin{theorem}[Rank-one local calculation]
Let \(\bar y\to Y\) be a geometric point.  For the origin
\(0\hookrightarrow\mathbf A^1_{\bar y}\), there is a canonical isomorphism

\[
 H^2_{\{0\}}(\mathbf A^1_{\bar y},\Lambda(1))\xrightarrow{\sim}\Lambda.
\]

Under this isomorphism, the boundary of the Kummer class of the coordinate
\(t\) is \(1\in\Lambda\).  All other supported cohomology groups of
\(\Lambda(1)\) vanish.
\end{theorem}

\begin{theorem}[Global existence, uniqueness, and base change of the orientation]
The Kummer boundary for the global coordinate \(t\) determines a morphism

\[
 \operatorname{pur}_s:\Lambda_Y
 \longrightarrow Rs^!\Lambda_{\mathbf A^1_Y}(1)[2].
\]

The target is a lisse rank-one \(\Lambda_Y\)-module concentrated in degree
zero after the displayed shift.  Restriction to geometric points gives an
injective map

\[
 \operatorname{Hom}_{D(Y,\Lambda)}
 \bigl(\Lambda_Y,Rs^!\Lambda_{\mathbf A^1_Y}(1)[2]\bigr)
 \longrightarrow
 \prod_{\bar y\to Y}
 H^2_{\{0\}}(\mathbf A^1_{\bar y},\Lambda(1)).
\]

The image contains the family whose every component is \(1\).
The morphism \(\operatorname{pur}_s\) is the unique global morphism with normalized fiber
class \(1\).  For every \(v:Y'\to Y\), the exceptional base-change
comparison identifies \(Lv^*\operatorname{pur}_s\) with the corresponding
orientation for the zero section of \(\mathbf A^1_{Y'}\).
\end{theorem}

\begin{definition}[Affine-line orientation class]
The affine-line orientation is the morphism supplied by the preceding
theorem,

\[
 \operatorname{pur}_s:\Lambda_Y
 \longrightarrow Rs^!\Lambda_{\mathbf A^1_Y}(1)[2]
\]

whose pullback to every geometric point of \(Y\) corresponds to
\(1\in\Lambda\) under the rank-one local calculation.
\end{definition}

\begin{proposition}[Support and base change of the orientation]
The adjoint class

\[
 s_*\Lambda_Y\longrightarrow
 R\Gamma_{s(Y)}\bigl(\Lambda_{\mathbf A^1_Y}(1)[2]\bigr)
\]

restricts to zero on \(\mathbf G_{m,Y}\).  For every morphism \(Y'\to Y\),
pullback identifies \(\operatorname{pur}_s\) with the orientation class of
the zero section of \(\mathbf A^1_{Y'}\).
\end{proposition}

\begin{theorem}[Affine-line integration morphism]
There is a unique natural transformation

\[
 \int_p:Rp_!\Lambda_{\mathbf A^1_Y}(1)[2]\longrightarrow\Lambda_Y
\]

whose pullback to every geometric point \(\bar y\to Y\) sends the compactly
supported fundamental class determined by \(\operatorname{pur}_s\) to
\(1\in\Lambda\).  It is compatible with arbitrary base change on \(Y\).
\end{theorem}

\begin{definition}[Affine-line purity comparison]
For \(K\in\Dctf(Y,\Lambda)\), the projection formula and affine-line
integration give the typed composite

\[
 Rp_!Lp^*K(1)[2]
 \xrightarrow{\sim}
 K\otimes^L_{\Lambda_Y}Rp_!\Lambda_{\mathbf A^1_Y}(1)[2]
 \xrightarrow{\operatorname{id}_K\otimes\int_p}K.
\]

Define

\[
 \gamma_{p,K}:Lp^*K(1)[2]\longrightarrow Rp^!K
\]

to be its adjoint under \(Rp_!\dashv Rp^!\).
\end{definition}

\begin{theorem}[Affine-line trace calculation]
The composite

\[
 Rp_!Lp^*K(1)[2]
 \xrightarrow{Rp_!(\gamma_{p,K})}
 Rp_!Rp^!K
 \xrightarrow{\Tr_{p,K}}K
\]

is the standard compactly supported integration map and sends the normalized
fiber orientation to \(1\in\Lambda\).  Equivalently, for every geometric
\(\bar y\to Y\), the induced map

\[
 H_c^2(\mathbf A^1_{\bar y},\Lambda(1))\xrightarrow{\sim}\Lambda
\]

is the isomorphism taking the fundamental class to \(1\).
\end{theorem}

\begin{theorem}[Affine-line purity]
For every \(K\in\Dctf(Y,\Lambda)\), the affine-line comparison is an
isomorphism

\[
 \gamma_{p,K}:Lp^*K(1)[2]\xrightarrow{\sim}Rp^!K.
\]

It is natural in \(K\), compatible with arbitrary base change on \(Y\), and
compatible with the exceptional trace.
\end{theorem}

\section{Affine space and etale localization}

\begin{definition}[Affine-space comparison]
Write \(p_d:\mathbf A^d_Y\to Y\) as the ordered composite of \(d\)
affine-line projections.  Define

\[
 \gamma_{p_d,K}:Lp_d^*K(d)[2d]\longrightarrow Rp_d^!K
\]

as the ordered composite of the \(d\) affine-line purity maps, the
composition isomorphisms for exceptional inverse image, and the additivity
isomorphisms for Tate twists.
\end{definition}

\begin{proposition}[Affine-space composition law]
Under
\(\mathbf A^{d+e}_Y\cong\mathbf A^d_{\mathbf A^e_Y}\), the equality

\[
 \gamma_{p_{d+e},K}
 =c_{p_e,p_d}^{-1}\circ
 \gamma_{p_d,Rp_e^!K}\circ
 Lp_d^*(\gamma_{p_e,K})(d)[2d]
\]

holds after the canonical reassociation of twists and shifts.
\end{proposition}

\begin{corollary}[Affine-space purity]
The map \(\gamma_{p_d,K}\) is an isomorphism for every
\(K\in\Dctf(Y,\Lambda)\), is symmetric under permutation of affine
coordinates, and is compatible with base change and products.
\end{corollary}

\begin{definition}[Etale coordinate chart]
Let \(f:X\to Y\) be smooth of pure relative dimension \(d\).  An etale
coordinate chart is an open \(X_i\subseteq X\) and an etale map

\[
 u_i:X_i\longrightarrow\mathbf A^d_Y
\]

such that \(f_i=f|_{X_i}=p_d u_i\).  The etale normalization
\(Ru_i^!\cong Lu_i^*\) and affine-space purity define

\[
 \gamma_{i,K}:Lf_i^*K(d)[2d]\longrightarrow Rf_i^!K.
\]
\end{definition}

\begin{proposition}[Etale change of coordinates]
Let \(u_i\) and \(u_j\) be two etale coordinate charts.  Their purity maps
have equal restrictions to \(X_i\cap X_j\).  The equality is stable under a
common etale refinement and is compatible with the composition constraint
for \(Rf^!\).
\end{proposition}

\begin{definition}[Mapping complex]
At the finite coefficient stage, put

\[
 \operatorname{Map}_X(M,N)=
 R\Gamma\bigl(X,R\mathcal Hom_X(M,N)\bigr).
\]

For the fixed replacements, degree-zero cocycles in
\(\mathcal Hom_X^\bullet(P(M),I(N))\) are cochain maps, and degree
\((-1)\) elements whose differential is the difference of two such
cocycles are cochain homotopies between them.
\end{definition}

\begin{proposition}[Mapping-complex model for derived morphisms]
The displayed differential squares to zero.  The resulting derived internal
Hom and mapping complex are independent of the selected replacements up to
canonical isomorphism, naturally in \(M\) and \(N\), and

\[
 \operatorname{Hom}_{D(X,\Lambda)}(M,N)
 \cong H^0\operatorname{Map}_X(M,N).
\]
\end{proposition}

\begin{definition}[The strict Cech mapping total complex]
For a finite ordered open cover \((X_i)\), set
\[
 C^{r,s}(M,N)=
 \prod_{i_0<\cdots<i_r}
 \Gamma\!\left(
 X_{i_0\ldots i_r},
 \mathcal Hom_X^s(P(M),I(N))|_{X_{i_0\ldots i_r}}
 \right).
\]

Let
\[
 \delta=\sum_{a=0}^{r+1}(-1)^a\operatorname{res}_a:
 C^{r,s}\longrightarrow C^{r+1,s}
\]
be the alternating Cech differential.  Define
\[
 \operatorname{Tot}^n C=
 \bigoplus_{r+s=n}C^{r,s},
 \qquad
 d_{\operatorname{Tot}}=\delta+(-1)^r d_{\mathcal Hom}
 \quad\hbox{on }C^{r,s}.
\]

The cover is finite, so every displayed total-degree sum is finite.
\end{definition}

\begin{proposition}[Cech total differential]
The alternating restriction operator and the internal-Hom differential
satisfy
\[
 \delta^2=0,\qquad d_{\mathcal Hom}^2=0,\qquad
 \delta d_{\mathcal Hom}=d_{\mathcal Hom}\delta.
\]
The equality \(d_{\operatorname{Tot}}^2=0\) holds, and
\(\operatorname{Tot}C^{\bullet,\bullet}(M,N)\) is an ordinary cochain complex.
\end{proposition}

\begin{theorem}[Cech--internal-Hom descent and strictification]
Let \((X_i\subseteq X)\) be a finite open cover, let
\(X_{i_0\ldots i_r}=X_{i_0}\cap\cdots\cap X_{i_r}\), and let
\(M,N\in\Dbc(X,\Lambda)\).  Define the Cech cosimplicial complex by

\[
 C^r(M,N)=
 \prod_{i_0,\ldots,i_r}
 R\Gamma\!\left(X_{i_0\ldots i_r},
 R\mathcal Hom(M,N)|_{X_{i_0\ldots i_r}}\right).
\]

The augmentation to its totalization is an isomorphism

\[
 R\Gamma\bigl(X,R\mathcal Hom_X(M,N)\bigr)
 \xrightarrow{\sim}\operatorname{Tot}C^\bullet(M,N).
\]

Restriction identifies global derived morphisms with local derived morphisms
\(u_i:M|_{X_i}\to N|_{X_i}\) equipped with equality homotopies on double
intersections and the coherent cocycle data on all higher intersections.  For
the functorial K-flat source, K-injective target, and their Cech total
resolutions, these data are represented by strict cocycles of cochain maps.
Pairwise equality in the derived category is not itself a coherent descent
datum.

A derived morphism is an isomorphism if its restriction to every \(X_i\) is
an isomorphism.
\end{theorem}

\begin{theorem}[Coherent descent datum for coordinate purity maps]
For every finite family of etale coordinate charts, the change-of-coordinate
identifications of the local maps \(\gamma_{i,K}\) extend from double
intersections to a Cech descent datum on all multiple intersections.  These
data are natural in \(K\), stable under refinement, and compatible with the
composition isomorphism for exceptional inverse image.
\end{theorem}

\begin{definition}[Glued smooth-purity map]
Choose a finite open cover \((X_i)\) equipped with etale coordinate maps
\(u_i:X_i\to\mathbf A^d_Y\).  The glued smooth-purity map is the unique
natural transformation supplied by Cech--internal-Hom descent from the
coherent coordinate-purity datum,

\[
 \gamma_{f,K}:Lf^*K(d)[2d]\longrightarrow Rf^!K
\]

whose restriction to \(X_i\) is \(\gamma_{i,K}\).
\end{definition}

\section{Smooth purity and coherence}

\begin{theorem}[Finite-coefficient smooth purity]
Let \(f:X\to Y\) be smooth and separated of pure relative dimension \(d\).
For every \(K\in\Dctf(Y,\Lambda)\), the glued comparison is an isomorphism

\[
 \gamma_{f,K}:Lf^*K(d)[2d]\xrightarrow{\sim}Rf^!K.
\]

It is independent of the etale-coordinate cover and is natural in \(K\).
\end{theorem}

\begin{theorem}[Composition coherence]
Let \(f:X\to Y\) and \(g:Y\to Z\) be smooth of pure relative dimensions
\(d\) and \(e\).  Under the canonical identification

\[
 L(gf)^*(-)(d+e)[2d+2e]
 \cong Lf^*Lg^*(-)(d)[2d](e)[2e],
\]

the map \(\gamma_{gf}\) equals the composite

\[
 L(gf)^*(-)(d+e)[2d+2e]
 \xrightarrow{Lf^*(\gamma_g)(d)[2d]}
 Lf^*Rg^!(-)(d)[2d]
 \xrightarrow{\gamma_f}
 Rf^!Rg^!(-)
 \xrightarrow{c_{g,f}^{-1}}R(gf)^!(-).
\]
\end{theorem}

\begin{definition}[Smooth pullback exchange for exceptional inverse image]
For a Cartesian square

\[
\begin{CD}
X' @>{u}>> X\\
@V{f'}VV @VV{f}V\\
Y' @>{v}>> Y
\end{CD}
\]

with \(f\) smooth of pure relative dimension \(d\), define

\[
 \operatorname{Ex}^{!,\mathrm{sm}}_{f,v,K}:
 Lu^*Rf^!K\longrightarrow Rf'^!Lv^*K
\]

by the typed composite

\[
\begin{aligned}
 Lu^*Rf^!K
 &\xrightarrow{(Lu^*\gamma_{f,K})^{-1}}
 Lu^*Lf^*K(d)[2d]\\
 &\xrightarrow{\sim}Lf'^*Lv^*K(d)[2d]
 \xrightarrow{\gamma_{f',Lv^*K}}Rf'^!Lv^*K.
\end{aligned}
\]

This is distinct from the right-adjoint mate
\(\operatorname{BC}^!_*:Rf^!Rv_*\to Ru_*Rf'^!\); both source and target
are fixed explicitly by their definitions.
\end{definition}

\begin{theorem}[Smooth base-change coherence]
For a Cartesian square

\[
\begin{CD}
X' @>{u}>> X\\
@V{f'}VV @VV{f}V\\
Y' @>{v}>> Y
\end{CD}
\]

with \(f\) smooth of pure relative dimension \(d\), the pullback of
\(\gamma_f\) is \(\gamma_{f'}\).  More precisely, the square

\[
\begin{CD}
Lu^*Lf^*K(d)[2d] @>{Lu^*\gamma_{f,K}}>> Lu^*Rf^!K\\
@V{\sim}VV @VV{\operatorname{Ex}^{!,\mathrm{sm}}_{f,v,K}}V\\
Lf'^*Lv^*K(d)[2d] @>{\gamma_{f',Lv^*K}}>> Rf'^!Lv^*K
\end{CD}
\]

commutes.  These squares are transitive under successive base changes.
\end{theorem}

\begin{corollary}[Tensor form of purity]
For \(K\in\Dctf(Y,\Lambda)\), there are coherent isomorphisms

\[
 Rf^!K\cong Lf^*K\otimes^L_{\Lambda_X}Rf^!\Lambda_Y
 \cong Lf^*K(d)[2d].
\]
\end{corollary}

\section{Trace and the structural morphism}

\begin{definition}[Smooth trace]
For a smooth separated \(f:X\to Y\) of pure relative dimension \(d\), define

\[
 \Tr_f^{\mathrm{sm}}:
 Rf_!\Lambda_X(d)[2d]
 \xrightarrow{Rf_!(\gamma_{f,\Lambda_Y})}
 Rf_!Rf^!\Lambda_Y
 \xrightarrow{\Tr_f}\Lambda_Y.
\]

For arbitrary \(L\in\Dctf(Y,\Lambda)\), define the coefficient-valued smooth
trace by

\[
 \Tr^{\mathrm{sm}}_{f,L}:
 Rf_!Lf^*L(d)[2d]
 \xrightarrow{Rf_!(\gamma_{f,L})}
 Rf_!Rf^!L
 \xrightarrow{\Tr_{f,L}}L.
\]

For \(L=\Lambda_Y\), this is the preceding normalized trace.  The projection
formula identifies it with
\(L\otimes^L\Tr_f^{\mathrm{sm}}\).
\end{definition}

\begin{theorem}[Trace normalization and transitivity]
On every geometric fiber, the smooth trace is the product of the normalized
affine-line traces.  It sends the ordered fundamental class of
\(\mathbf A^d\) to \(1\).  For composable smooth maps \(f\) and \(g\),

For \(K\in\Dctf(Z,\Lambda)\), the typed composite

\[
\begin{aligned}
R(gf)_!L(gf)^*K(d+e)[2d+2e]
&\xrightarrow{\alpha_{g,f}^{-1}}
Rg_!Rf_!Lf^*Lg^*K(d+e)[2d+2e]\\
&\xrightarrow{Rg_!(\Tr^{\mathrm{sm}}_{f,Lg^*K(e)[2e]})}
Rg_!Lg^*K(e)[2e]\\
&\xrightarrow{\Tr^{\mathrm{sm}}_{g,K}}K
\end{aligned}
\]

equals \(\Tr^{\mathrm{sm}}_{gf,K}\), after the pullback-composition and
twist-additivity identifications.  Thus the \(f\)-trace, with the
\((e)[2e]\) twist already present, is applied before the \(g\)-trace.
\end{theorem}

\begin{corollary}[Finite-stage smooth dualizing complex]
Let \(A=\Lambda\) be a finite coefficient ring, let \(k\) be separably
closed, and let \(a_X:X\to\operatorname{Spec}k\) be separated, of finite
type, smooth, and of pure dimension \(d\).  The dualizing complex

\[
 \omega_X=Ra_X^!A
\]

has the canonical identification

\[
 \omega_X\xrightarrow{\sim}A_X(d)[2d],
\]

and the counit \(Ra_{X!}\omega_X\to A\) is the normalized integration map

\[
 R\Gamma_c(X,A_X(d)[2d])\longrightarrow A.
\]
\end{corollary}

\section{The fundamental class of a regular immersion}

Let \(k\) be a perfect field, \(X\) smooth over \(k\), and
\(i:Z\hookrightarrow X\) a regular closed immersion of pure codimension
\(c\) with \(Z\) regular.  Cohomology with supports is computed with
\(Ri_!Ri^!=i_*Ri^!\), and a class of degree \(n\) supported on \(Z\) with
coefficients \(L\) is a morphism \(\Lambda_X\to Ri_!Ri^!(L[n])\).

\begin{definition}[Kummer boundary]
Let \(f\in\Gamma(U,\mathcal O_U)\), let \(i_f:V(f)\hookrightarrow U\) and
let \(j_f:U_f\hookrightarrow U\) be its open complement, on which \(f\) is a
unit.  The Kummer boundary

\[
 \partial\kappa(f)\in H^2_{V(f)}(U,\Lambda(1))
\]

is the image of the Kummer class
\(\kappa(f|_{U_f})\in H^1(U_f,\Lambda(1))\), read as a class
\(\Lambda_U\to Rj_{f*}Lj_f^*\Lambda_U(1)[1]\) by adjunction, under the
boundary of the localization triangle of the pair \((i_f,j_f)\).
\end{definition}

\begin{definition}[Cup product with supports]
For closed immersions \(i_1:Z_1\hookrightarrow U\),
\(i_2:Z_2\hookrightarrow U\) and \(i:Z_1\cap Z_2\hookrightarrow U\), the
tensor product of the counits
\(Ri_{1!}Ri_1^!M\otimes Ri_{2!}Ri_2^!N\to M\otimes N\) factors uniquely
through the counit \(Ri_!Ri^!(M\otimes N)\to M\otimes N\), because its source
is supported on \(Z_1\cap Z_2\).  Composing with the diagonal of
\(\Lambda_U\) defines

\[
 H^a_{Z_1}(U,M)\times H^b_{Z_2}(U,N)\longrightarrow
 H^{a+b}_{Z_1\cap Z_2}(U,M\otimes N),\qquad
 (\alpha,\beta)\longmapsto\alpha\cup\beta.
\]

Similarly, for a morphism \(e:U\to X\) and the base change
\(i':Z'\hookrightarrow U\) of a closed immersion \(i:Z\hookrightarrow X\),
the composite \(Le^*Ri_!Ri^!L\to Le^*L\) factors uniquely through
\(Ri'_!Ri'^!Le^*L\); this defines the restriction of supports
\(H^n_Z(X,L)\to H^n_{Z'}(U,Le^*L)\).
\end{definition}

\begin{definition}[Frames and frame classes]
A frame of \(i\) is an etale map \(e:U\to X\) together with functions
\(f_1,\dots,f_c\in\Gamma(U,\mathcal O_U)\) such that
\(V(f_1,\dots,f_c)=Z\times_XU\) as closed subschemes of \(U\).  Its frame
class is

\[
 \operatorname{cl}(f)=\partial\kappa(f_1)\cup\cdots\cup\partial\kappa(f_c)
 \in H^{2c}_{Z\times_XU}(U,\Lambda(c)),
\]

combined with \(\Lambda(1)[2]^{\otimes c}\cong\Lambda(c)[2c]\).  The empty
product is \(1\in H^0_U(U,\Lambda)\).  Frames exist Zariski-locally on \(X\)
because \(i\) is a regular immersion.
\end{definition}

\begin{theorem}[Fundamental class of a regular immersion]
\begin{enumerate}
\item The complex \(Ri^!\Lambda_X(c)[2c]\) is concentrated in degree zero,
  and its \(\mathcal H^0\) is locally free of rank one.
\item There is a unique class \([Z\subset X]\in H^{2c}_Z(X,\Lambda(c))\)
  whose restriction to every frame is the frame class.  The induced map
  \(\operatorname{cl}_i:\Lambda_Z\to Ri^!\Lambda_X(c)[2c]\),
  \(1\mapsto[Z\subset X]\), is an isomorphism.
\item The class is compatible with etale base change,
  \([Z\times_XX'\subset X']=[Z\subset X]|_{X'}\), and with exterior products,
  \([Z_1\times_kZ_2\subset X_1\times_kX_2]=p_1^*[Z_1\subset X_1]\cup
  p_2^*[Z_2\subset X_2]\).
\end{enumerate}
\end{theorem}

\begin{proof}[Sketch]
Part (1) follows from smooth purity for \(X\) and \(Z\) over \(k\) and the
co-projection formula for \(Ri^!\).  Two frames on the same \(U\) differ
Zariski-locally by a matrix in \(\mathrm{GL}_c(\mathcal O)\).  Such a matrix
is a product of diagonal, permutation and elementary matrices.  The frame
class is invariant under each of these: under diagonal matrices because
\(\partial\kappa(u)=0\) for a unit \(u\) and \(\kappa\) is additive; under
permutations because all degrees are even; and under elementary matrices
by \(\mathbf A^1\)-invariance.  So the frame classes depend only on
\(Z\times_XU\subset U\).  By (1) they are sections of the sheaf
\(\mathcal H^0(Ri^!\Lambda_X(c)[2c])\) over an etale cover, compatible on
overlaps, and they glue to a unique global class.  The rank-one local
calculation, applied to a frame, shows that \([Z\subset X]\) generates
\(\mathcal H^0\) at every geometric point, which gives (2).  Etale maps
pull frames back to frames, and the product of frames on \(U_1\) and
\(U_2\) is a frame on \(U_1\times_kU_2\), which gives (3).
\end{proof}

\section{Transport through coefficient stages}

\begin{theorem}[Finite reduction of purity]
For \(m\leq n\), reduction carries every ingredient of
\(\gamma_{f,\Lambda_n}\) to the corresponding ingredient over
\(\Lambda_m\).  The square

\[
\begin{CD}
\red_{n,m,X}(Lf^*K(d)[2d]) @>{\red(\gamma_{f,K})}>>
\red_{n,m,X}(Rf^!_{\Lambda_n}K)\\
@V{\sim}VV @VV{\sim}V\\
Lf^*\red_{n,m,Y}K(d)[2d] @>{\gamma_{f,\red K}}>>
Rf^!_{\Lambda_m}\red_{n,m,Y}K
\end{CD}
\]

commutes.  The same is true for affine orientations, traces, composition,
and base-change squares.
\end{theorem}

\begin{theorem}[Adic reconstruction]
For \(K=(K_n)\in\Dctf(Y,\OO_E)\), the maps
\(\gamma_{f,K_n}\) form a morphism of normalized systems.  They induce an
isomorphism

\[
 Lf^*K(d)[2d]\xrightarrow{\sim}Rf^!K
\]

in \(\Dctf(X,\OO_E)\).  Its reduction at every level is exactly
\(\gamma_{f,K_n}\).
\end{theorem}

\begin{theorem}[Rational and \(\Qlb\) purity]
Rationalization and finite extension of coefficient fields carry the adic
purity isomorphism to the purity isomorphism over \(E\).  These isomorphisms
are transitive in coefficient extensions and preserve idempotent summands.
They define

\[
 Lf^*K(d)[2d]\xrightarrow{\sim}Rf^!K
\]

for \(A=E\) and \(A=\Qlb\), with the same trace, composition, and base-change
coherence.
\end{theorem}

\begin{corollary}[Adic and rational smooth dualizing complexes]
Let \(A=\OO_E,E\), or \(\Qlb\), let \(k\) be separably closed, and let
\(a_X:X\to\operatorname{Spec}k\) be separated, of finite type, smooth, and
of pure dimension \(d\).  Coefficient transport gives

\[
 \omega_{X,A}=Ra_X^!A\xrightarrow{\sim}A_X(d)[2d],
\]

and identifies the counit with the transported normalized integration map

\[
 R\Gamma_c(X,A_X(d)[2d])\longrightarrow A.
\]
\end{corollary}

\begin{definition}[Scheme and geometric Frobenius]
Let \(f_0:X_0\to Y_0\) be smooth and separated of pure relative dimension
\(d\) over \(k_0=\mathbf F_q\).  Put \(\bar k=\overline{\mathbf F}_q\),
\(X=X_0\times_{k_0}\bar k\), \(Y=Y_0\times_{k_0}\bar k\), and let
\(f:X\to Y\) be the base change of \(f_0\).  Write

\[
 \sigma(z)=z^q,\qquad \gamma=\sigma^{-1},
\]

and denote the induced semilinear automorphisms by
\(g_{X,\gamma}:X\to X\) and \(g_{Y,\gamma}:Y\to Y\).  They form the
Cartesian square

\[
\begin{CD}
X @>{g_{X,\gamma}}>> X\\
@V{f}VV @VV{f}V\\
Y @>{g_{Y,\gamma}}>> Y.
\end{CD}
\]

Let \(\Phi_X\) and \(\Phi_Y\) be the base changes to \(\bar k\) of the
\(q\)-power scheme Frobenius maps of \(X_0\) and \(Y_0\).  Topological
invariance supplies canonical natural isomorphisms

\[
 \kappa_X:L\Phi_X^*\xrightarrow{\sim}Lg_{X,\gamma}^*,
 \qquad
 \kappa_Y:L\Phi_Y^*\xrightarrow{\sim}Lg_{Y,\gamma}^*.
\]

Under the pullback-composition identifications, these satisfy
\(Lf^*\kappa_Y=\kappa_XLf^*\).
\end{definition}

\begin{definition}[Strict Weil data and the Tate datum]
At a finite coefficient stage \(A=\Lambda\), let
\(K\in\Dctf(Y,\Lambda)\).  A strict geometric-Weil datum on \(K\) is a
specified isomorphism in the ordinary triangulated category

\[
 \varphi_K:Lg_{Y,\gamma}^*K\xrightarrow{\sim}K.
\]

For \(A=\OO_E\), such a datum means a normalized system
\((K_n,\varphi_{K_n})_{n\geq1}\) of finite-stage data: the reduction
isomorphisms
\[
 \red_{n+1,n,Y}K_{n+1}\xrightarrow{\sim}K_n
\]
intertwine \(\red_{n+1,n,Y}(\varphi_{K_{n+1}})\) with
\(\varphi_{K_n}\), and the underlying complexes have common bounds,
stratification, and Tor-amplitude.  For \(A=E\), a datum is an object obtained
from an \(\OO_E\)-datum by inverting \(\varpi\), together with an idempotent
summand when required.  For \(A=\Qlb\), it is represented by an
idempotent-equipped datum over some finite \(E/\Ql\), modulo further finite
coefficient extension.  In every case the forgetful object lies in
\(\Dctf(Y,A)\), and the symbol \(\varphi_K\) denotes the resulting typed
isomorphism \(Lg_{Y,\gamma}^*K\xrightarrow{\sim}K\).

Its scheme-Frobenius form is the typed composite

\[
 u_K=\varphi_K\circ\kappa_{Y,K}:
 L\Phi_Y^*K\longrightarrow K.
\]

The Tate line has its canonical datum
\(\tau_{X,1}:Lg_{X,\gamma}^*A_X(1)\xrightarrow{\sim}A_X(1)\), on which
geometric Frobenius acts by \(q^{-1}\).  Tensor powers and tensor inverses
give

\[
 \tau_{X,r}:Lg_{X,\gamma}^*A_X(r)\xrightarrow{\sim}A_X(r),
\]

with geometric-Frobenius scalar \(q^{-r}\).  The induced datum on the
source of smooth purity is

\[
\begin{aligned}
 \varphi^{\mathrm{pur}}_{f,K}:
 Lg_{X,\gamma}^*\bigl(Lf^*K(d)[2d]\bigr)
 &\xrightarrow{\sim}
 Lf^*Lg_{Y,\gamma}^*K\otimes_{A_X}^L
 Lg_{X,\gamma}^*A_X(d)[2d]\\
 &\xrightarrow{Lf^*\varphi_K\otimes\tau_{X,d}}
 Lf^*K(d)[2d].
\end{aligned}
\]

The datum on the exceptional inverse image is

\[
\begin{aligned}
 \varphi^!_{f,K}:Lg_{X,\gamma}^*Rf^!K
 &\xrightarrow{\operatorname{Ex}^{!,\mathrm{sm}}_{f,g_{Y,\gamma},K}}
 Rf^!Lg_{Y,\gamma}^*K
 \xrightarrow{Rf^!(\varphi_K)}Rf^!K.
\end{aligned}
\]
\end{definition}

\begin{definition}[Compact-support transport of a Weil datum]
For \(M\in\Dctf(X,A)\) with a specified datum
\(\varphi_M:Lg_{X,\gamma}^*M\xrightarrow{\sim}M\), define

\[
\begin{aligned}
 \varphi_{f_!M}:Lg_{Y,\gamma}^*Rf_!M
 &\xrightarrow{\operatorname{BC}_!}
 Rf_!Lg_{X,\gamma}^*M
 \xrightarrow{Rf_!(\varphi_M)}Rf_!M.
\end{aligned}
\]

For every object \(N\) on \(X\) or \(Y\) carrying a datum \(\varphi_N\),
its scheme-Frobenius form is

\[
 u_N=\varphi_N\circ\kappa_{T,N}:L\Phi_T^*N\longrightarrow N,
\]

where \(T=X\) or \(Y\) is the base of \(N\), and \(\kappa_{T,N}\) is the
component of \(\kappa_T\) at \(N\).
\end{definition}

\begin{theorem}[Weil equivariance of purity and trace]
For every strict Weil datum \((K,\varphi_K)\), smooth purity is equivariant:

\[
\begin{CD}
Lg_{X,\gamma}^*Lf^*K(d)[2d]
 @>{Lg_{X,\gamma}^*(\gamma_{f,K})}>>
Lg_{X,\gamma}^*Rf^!K\\
@V{\varphi^{\mathrm{pur}}_{f,K}}V{\sim}V
@VV{\varphi^!_{f,K}}V\\
Lf^*K(d)[2d] @>{\gamma_{f,K}}>> Rf^!K.
\end{CD}
\]

If \(\varphi_{f_!Rf^!K}\) is formed from \(\varphi^!_{f,K}\), then the
exceptional counit is equivariant:

\[
\begin{CD}
Lg_{Y,\gamma}^*Rf_!Rf^!K
 @>{Lg_{Y,\gamma}^*(\Tr_{f,K})}>> Lg_{Y,\gamma}^*K\\
@V{\varphi_{f_!Rf^!K}}V{\sim}V @VV{\varphi_K}V\\
Rf_!Rf^!K @>{\Tr_{f,K}}>> K.
\end{CD}
\]

If \(\varphi_{f_!Lf^*K(d)[2d]}\) is formed from
\(\varphi^{\mathrm{pur}}_{f,K}\), then the normalized smooth trace is
equivariant:

\[
\begin{CD}
Lg_{Y,\gamma}^*Rf_!Lf^*K(d)[2d]
 @>{Lg_{Y,\gamma}^*(\Tr^{\mathrm{sm}}_{f,K})}>>
Lg_{Y,\gamma}^*K\\
@V{\varphi_{f_!Lf^*K(d)[2d]}}V{\sim}V @VV{\varphi_K}V\\
Rf_!Lf^*K(d)[2d] @>{\Tr^{\mathrm{sm}}_{f,K}}>> K.
\end{CD}
\]

Precomposing every vertical map with the corresponding \(\kappa_X\) or
\(\kappa_Y\) gives the same three squares for the scheme-Frobenius maps
\(u_N\).  All geometric- and scheme-Frobenius squares commute with finite
reduction, scalar extension, composition, and base change.
\end{theorem}

\end{document}
